\documentclass[10pt,twocolumn]{article}
\usepackage[a4paper,margin=1.8cm]{geometry}
\usepackage{booktabs}
\usepackage{array}
\usepackage{url}
\usepackage[hidelinks]{hyperref}
\usepackage{enumitem}
\setlist{nosep,leftmargin=*}

\newcommand{\code}[1]{\texttt{\small #1}}
\renewcommand{\arraystretch}{1.25}

\title{\vspace{-1.2em}RFC 10070 client-authentication audience rule:\\
a five-implementation conformance snapshot}
\author{Nobuo Otoi\\[2pt]
\small ba-cybersecurity.com, Tokyo, Japan\\
\small independent cybersecurity researcher}
\date{\small Draft note, 2026-11-09}

\begin{document}
\twocolumn[
\begin{@twocolumnfalse}
\maketitle
\begin{abstract}
\noindent RFC 10070, an IETF standard published 2026-11, tightened the audience
(\code{aud}) rule for the JWT used in OAuth client authentication: the \code{aud} MUST be the authorization
server's issuer identifier as the sole value, and the token-endpoint URL MUST NOT
be used. A read-only source review of five major open OAuth authorization-server
implementations finds that, against their default-branch HEADs, \emph{all five
still accept the token-endpoint URL as the client-authentication \code{aud}} --- the
value RFC 10070 now forbids --- and four also accept a multi-valued \code{aud}. The
per-implementation severity is Low (conformance, not a new flaw): the underlying
issue is a previously disclosed, spec-level vulnerability (OpenID Foundation,
2025-01; ePrint 2025/629), already fixed by RFC 10070, for which server-side
acceptance of the endpoint URL is only a necessary condition. This is a timed
measurement of ecosystem adoption of a published fix, not a scorecard.
\vspace{0.8em}
\end{abstract}
\end{@twocolumnfalse}
]

\begin{table*}[t]
\centering
\footnotesize
\caption{Client-authentication \code{aud} handling at default-branch HEAD, 2026-11-09. ``Divergent'' = accepts the token-endpoint URL that RFC 10070 now forbids; ``mixed'' = partially aligned.}
\label{tab:matrix}
\begin{tabular}{@{}p{2.2cm}p{5.0cm}p{2.1cm}p{3.6cm}p{1.6cm}@{}}
\toprule
\textbf{Implementation} & \textbf{Accepted \code{aud} (client auth)} & \textbf{Multi-valued \code{aud}} & \textbf{Conformant by config alone?} & \textbf{Class} \\
\midrule
ory/fosite (in Hydra) & configured token URL(s) only; issuer identifier \emph{not} accepted & accepted & partial (no clean path) & divergent \\
node-oidc-provider & issuer / token / called-endpoint URL (default); issuer-only under FAPI~2.0 & accepted (default); rejected under FAPI~2.0 & yes (hook; strict under FAPI~2.0) & mixed \\
Keycloak & issuer / token / introspect / PAR / CIBA URLs & rejected by default & no (accept-set hardcoded) & mixed \\
Spring Authz Server & issuer / token / introspect / revoke / PAR URLs & accepted & yes (\code{setJwtValidatorFactory}) & divergent \\
Authlib & developer-configured; sole-value not enforced & accepted & partial (needs \code{verify\_claims} override) & divergent \\
\bottomrule
\end{tabular}
\end{table*}

\section{Summary}
RFC 10070 (``Updates to OAuth 2.0 JWT Client Authentication and Assertion-Based
Authorization Grants''; updates RFC 7521/7522/7523/9126) tightened the audience
rule for the JWT used in \textbf{client authentication} (\code{private\_key\_jwt} /
\code{client\_secret\_jwt}, and the PAR endpoint): the \code{aud} MUST be the
authorization server's \textbf{issuer identifier as the sole value}, and the
token-endpoint URL \textbf{MUST NOT} be used. The \textbf{authorization-grant}
path (JWT bearer, RFC 7523 \S3 item~3 case~A) stays lenient and is out of scope.

Against their default-branch HEADs as of 2026-11-09, \textbf{all five implementations
still accept the token-endpoint URL as the client-authentication \code{aud}}. Four
also accept a multi-valued \code{aud}. Two are partially aligned: Keycloak rejects
a multi-valued \code{aud} by default, and node-oidc-provider is strict under its
FAPI~2.0 profile.

\section{Why this is Low, not a vulnerability}
The per-implementation severity is \textbf{Low}, and this is not a report of a
novel flaw in any one project. But the underlying issue is real. The audience
restriction fixes a \textbf{previously disclosed, spec-level vulnerability}: a
malicious authorization server can advertise a legitimate server's token-endpoint
URL in its own metadata --- a client can authenticate a server's \emph{issuer}
identifier but not its endpoints (RFC 8414), as RFC 10070 \S1 notes --- so a client
that reuses one key pair across servers and sets \code{aud} to the endpoint URL
hands the malicious server a client assertion it can replay to the legitimate one.
This was disclosed across standards bodies (OpenID Foundation, 2025-01-24) and
analyzed in the literature (ePrint 2025/629); RFC 10070 \S1 calls it a ``security
vulnerability,'' and RFC 10070 is the fix.

A server accepting the token-endpoint URL as \code{aud} is a \textbf{necessary
condition} for that attack; the residual real-world risk depends on whether such
key-reusing, endpoint-URL-audience clients exist against a given server. So this
note measures \textbf{non-adoption of a published fix}, not a latent bug in any
single project --- hence Low, conformance-framed, and suitable (if raised at all)
as a conformance heads-up, not a vulnerability report.

\section{The rule is not new; adoption is phased}
The restriction has been public for roughly two years: disclosed 2025-01 (OIDF),
analyzed in 2025 (ePrint 2025/629), and carried through
\mbox{draft-ietf-oauth-rfc7523bis}, FAPI~2.0, and OIDC Core errata before RFC 10070
finalized it in 2026-11. Implementers have had notice and move on their own
schedules: Keycloak's 26.2 upgrade notes already recommend the issuer identifier as
\code{aud}, and Spring Security \#18226 records the maintainer deferring the change
until the draft was approved. RFC 10070 itself advises phased adoption --- \S6
recommends updating clients before or alongside server enforcement, and \S4 cautions
that rejecting JWTs without an explicit \code{typ} is NOT RECOMMENDED. A server
still accepting the old \code{aud} may be mid-transition by design. The value of
this snapshot is as a timed measurement of that transition.

\section{The rule (verified against RFC 10070)}
\begin{itemize}
\item Client authentication --- RFC 7523 \S3 item~3 case~B, as replaced by RFC 10070
  \S4: the \code{aud} ``MUST indicate only the issuer identifier of the intended
  authorization server''; the server ``MUST reject any JWT that does not contain its
  issuer identifier as its sole audience value''; the token-endpoint URL ``MUST NOT
  be used as an audience value.''
\item PAR --- RFC 9126 \S2, as replaced by RFC 10070 \S5: the issuer identifier
  ``MUST be used as the sole value of the audience''; reject otherwise.
\item Authorization grant (out of scope) --- RFC 7523 \S3 item~3 case~A: an
  authorization server ``MAY be identified by either its issuer identifier or its
  token endpoint URL.''
\item RFC 10070 does not forbid a single-element array; ``sole audience value'' is
  the operative phrase. No backward-compatibility text permits servers to keep
  accepting the old values.
\end{itemize}

\section{Method}
Read-only source review; no build, no network, nothing sent. For each project the
\textbf{client-authentication} \code{aud} validator was located and its accept-set
read at the pinned commit, deliberately distinguished from the authorization-grant
validator (fosite \code{handler/rfc7523/handler.go}, Keycloak
\code{JWTAuthorizationGrantType.java}, Authlib \code{jwt\_bearer.py}; node-oidc-provider
and Spring have no JWT-bearer grant). The multi-valued-\code{aud} behavior was
confirmed in code for each. The RFC 10070 normative text, the OIDF disclosure, and
the ePrint abstract were read against primary sources. Deciding code at the pinned
commits: fosite \code{client\_authentication.go} + \code{fositex/config.go:230--236}
(Hydra in-tree); node-oidc-provider \code{jwt\_client\_auth.js:48--58},
\code{oidc\_context.js:133--135}, \code{defaults.js:632--638}; Keycloak
\code{AbstractBaseJWTValidator.java:126--137}, \code{JWTClientValidator.java:29--39};
Spring \code{JwtClientAssertionDecoderFactory.java:174,194--212}; Authlib
\code{rfc7523/client.py:64--94}. Pins: fosite \code{a5f0b09} / Hydra \code{4174065} /
nop \code{3ed87be} / Keycloak \code{d4d5554} / Spring \code{4283973} / Authlib
\code{0f9826c}.

\section*{AI-assistance disclosure}
Most of the source review and drafting was performed with AI assistance. Before
finalizing, the author re-opened each cited file:line at its pinned commit and
confirmed the classification, and verified the RFC 10070 rules, the
client-auth-vs-grant split, and the disclosure history against the primary
documents.

\end{document}
